\section{Integral topology without the cusp-collapse argument}
\label{sec:integral-leray}

This section follows the independent nearby-cycle/Leray route.  It does not
use the candidate's cellular collapse, its Mayer--Vietoris homology matrix,
or the asserted equality between those two constructions.  All coefficient
groups are integral.

Put
\[
 B=\mathbb P^1,\quad \Sigma=\{p_0,p_1,p_2\},\quad
 B^\circ=B\setminus\Sigma,\quad j:B^\circ\hookrightarrow B,
\]
and let \(F\) be a smooth fibre.  Retain
\[
 V=H^1(F)=\mathbb Z\langle\gamma,u,w,\delta\rangle,
 \quad q=uw+6\gamma\delta,
 \quad \operatorname{vol}=\gamma uw\delta.
\]

\subsection{Integral specialization at the cusp}

\begin{theorem}[nearby-cycle specialization]
\label{thm:p0-integral-specialization}
\Verified\ For the proper holomorphic cusp filling
\(f_0:N_0\to\Delta\) of \cref{thm:candidate-cusp-geometry}, with its
proved local normal-crossings equations and normalization combinatorics,
the specialization map of the central fibre \(W\),
\[
 \operatorname{sp}^r:H^r(W;\mathbb Z)\longrightarrow H^r(F;\mathbb Z)
\]
is injective and
\[
 \operatorname{im}\operatorname{sp}^r
 =H^r(F;\mathbb Z)^{T_0}
\tag{5.1}
\]
for \(0\le r\le4\).  In particular, the image is saturated and there is no
unrecorded finite-index defect at \(p_0\).
\end{theorem}

\begin{proof}
Let \(\widetilde\Delta^*\to\Delta^*\) be the universal cover, put
\(\widetilde N_0^*=N_0\times_\Delta\widetilde\Delta^*\), and let
\[
 k:\widetilde N_0^*\longrightarrow N_0,
 \qquad i:W\hookrightarrow N_0
\]
be the natural maps.  The analytic nearby-cycle complex used here is
\[
 \psi=\psi_{f_0}\mathbb Z
 :=i^*Rk_*k^*\mathbb Z_{N_0}\in D_c^b(W).
\tag{5.0}
\]
Dimca defines analytic nearby cycles in
\cite[Definition~4.2.1]{Dimca2004}; the Milnor-fibre description of their
stalks is \cite[Proposition~4.2.2]{Dimca2004}.  The proper-direct-image
theorem for nearby cycles is
\cite[Proposition~4.2.11]{Dimca2004}; applied to \(f_0\) and the identity
map of the disc, it gives
\[
 \mathbb H^n(W;\psi)\simeq H^n(F;\mathbb Z)
\]
and identifies the morphism induced by the unit
\(\mathbb Z_W\to\psi\) with specialization.  Proper base change itself is
\cite[Theorem~2.3.26]{Dimca2004}.  Dimca's coefficient hypothesis is a
commutative noetherian ring of finite global dimension, which
\(\mathbb Z\) satisfies.  Thus the general integral analytic formalism is
available here.  The candidate-specific existence and local description of
\(f_0\) are proved in \cref{thm:candidate-cusp-geometry,% 
prop:candidate-cusp-overlap}.  Kashiwara--Schapira
\cite[\S8.6]{KashiwaraSchapira1990}
corroborates the analytic formalism, while SGA~7, Expos\'e~XIII
\cite{SGA7II1973}, is its algebraic antecedent.

At a point where \(f=z_0\cdots z_k\), \(k=0,1,2\), the Milnor fibre
deformation retracts onto \((S^1)^k\): in polar coordinates the radial
variables form a contractible open simplex and the phases satisfy one linear
relation.  Hence
\[
 \mathcal H^0\psi=\mathbb Z_W,\qquad
 \mathcal H^1\psi=:\mathscr V,\qquad
 \mathcal H^2\psi=\mathbb Z_P\oplus\mathbb Z_Q,\qquad
 \mathcal H^b\psi=0\ (b\ge3),
\tag{5.2}
\]
where \(\mathscr V\) is supported on the three double curves
\(\bar C_1,\bar C_2,\bar C_3\simeq\mathbb P^1\).

Orient each vanishing circle by the displayed order of its two branches in
the normalization.  At \(P\), the branches of
\((\bar C_1,\bar C_2,\bar C_3)\) are respectively
\[
 (b_0,b_1),\qquad(b_0,b_2),\qquad(b_1,b_2).
\]
The three restrictions are therefore evaluations on
\[
 e_0-e_1,\qquad e_0-e_2,\qquad e_1-e_2,
\]
whose alternating sum is zero.  At \(Q\), the ordered branch pairs are
\[
 (c_2,c_1),\qquad(c_0,c_1),\qquad(c_0,c_2),
\]
so the evaluations are
\[
 e_2-e_1,\qquad e_0-e_1,\qquad e_0-e_2,
\]
and again the first minus the second plus the third is zero.  The global
normalization labels \(C_k\) and \(C_{k+3}\) distinguish the two branches
of each double curve and therefore fix these signs simultaneously.  Thus
\[
0\longrightarrow\mathscr V
\longrightarrow\bigoplus_{k=1}^3 i_{k*}\mathbb Z_{\bar C_k}
\xrightarrow{\delta_1}\mathbb Z_P\oplus\mathbb Z_Q
\longrightarrow0,
\]
\[
 \delta_1(n_1,n_2,n_3)=(s,s),\qquad s=n_1-n_2+n_3.
\tag{5.3}
\]
The image \(\mathbb Z(1,1)\) is primitive.  Taking cohomology of (5.3)
gives free groups with ranks
\[
 \operatorname{rk}H^a(W;\mathscr V)=(2,1,3,0)
 \quad(a=0,1,2,3).
\tag{5.4}
\]

The normalization resolution for the constant sheaf is
\[
0\longrightarrow\mathbb Z_W
\longrightarrow\nu_*\mathbb Z_{\mathrm{dP}_6}
\longrightarrow\bigoplus_{k=1}^3 i_{k*}\mathbb Z_{\bar C_k}
\xrightarrow{\delta_1}\mathbb Z_P\oplus\mathbb Z_Q
\longrightarrow0.
\tag{5.5}
\]
This is exact stalkwise.  At a smooth point it is the identity
\(\mathbb Z\to\mathbb Z\).  At a generic double point it is the augmented
complex
\[
 0\longrightarrow\mathbb Z\longrightarrow\mathbb Z^2
 \xrightarrow{(x_0,x_1)\mapsto x_0-x_1}\mathbb Z\longrightarrow0.
\]
At a triple point, in either of the branch orders just fixed, it is
\[
 0\longrightarrow\mathbb Z\longrightarrow\mathbb Z^3
 \xrightarrow{d_0}\mathbb Z^3
 \xrightarrow{d_1}\mathbb Z\longrightarrow0,
\]
where at \(P\)
\[
 d_0(x_0,x_1,x_2)=(x_0-x_1,x_0-x_2,x_1-x_2),
 \qquad d_1(n_1,n_2,n_3)=n_1-n_2+n_3,
\]
while at \(Q\)
\[
 d_0(x_0,x_1,x_2)=(x_2-x_1,x_0-x_1,x_0-x_2),
 \qquad d_1(n_1,n_2,n_3)=n_1-n_2+n_3.
\]
These augmented simplicial complexes are exact, which proves (5.5),
including its endpoint signs.

On \(H^2\), the only nontrivial map is
\[
 d:H^2(\mathrm{dP}_6;\mathbb Z)\longrightarrow\mathbb Z^3,
 \qquad
 d(L)=(L\cdot C_k-L\cdot C_{k+3})_{k=1}^3.
\]
Write \(\mathrm{dP}_6\) as the blow-up of \(\mathbb P^2\) at three
noncollinear points, with basis \(H,E_1,E_2,E_3\), and order its hexagon as
\[
 E_1, H-E_1-E_2, E_2, H-E_2-E_3, E_3, H-E_1-E_3.
\]
For \(L=aH-\sum b_iE_i\), direct intersection gives
\[
 d(L)=(s,-s,s),\qquad s=b_1+b_2+b_3-a.
\tag{5.6}
\]
Thus \(\operatorname{im}d=\mathbb Z(1,-1,1)\) is primitive,
\(\ker d\simeq\mathbb Z^3\), and
\(\operatorname{coker}d\simeq\mathbb Z^2\).  The resolution (5.5) has
width two and the cohomology of its smooth strata is concentrated in even
degrees.  The only possible higher differential would be
\(d_2:E_2^{0,q}\to E_2^{2,q-1}\), whose target is an odd row and is zero;
differentials with \(r\ge3\) leave the width-two range.  Therefore
\[
 H^a(W;\mathbb Z)=
 (\mathbb Z,\mathbb Z^2,\mathbb Z^4,\mathbb Z^2,\mathbb Z)
 \quad(a=0,1,2,3,4),
\tag{5.7}
\]
with no torsion.

The hypercohomology spectral sequence
\[
 E_2^{a,b}=H^a(W;\mathcal H^b\psi)
 \Longrightarrow H^{a+b}(F;\mathbb Z)
\tag{5.8}
\]
has the following ranks, all on free groups:
\[
\begin{array}{c|ccccc}
b=2&2&0&0&0&0\\
b=1&2&1&3&0&0\\
b=0&1&2&4&2&1\\ \hline
&a=0&a=1&a=2&a=3&a=4.
\end{array}
\]
Let the possible differential ranks be
\[
\alpha=\operatorname{rk}d_2^{0,1},\quad
\beta=\operatorname{rk}d_2^{1,1},\quad
\gamma'=\operatorname{rk}d_2^{2,1},\quad
\delta'=\operatorname{rk}d_2^{0,2},\quad
\epsilon=\operatorname{rk}d_3^{0,2}.
\]
Comparison with the torus Betti numbers \((1,4,6,4,1)\) in total degrees
one, four, and two gives
\[
 \alpha=0,\qquad \gamma'=0,\qquad
 \beta+\delta'+\epsilon=1.
\tag{5.9}
\]

We determine the remaining ambiguity geometrically, keeping the supported
specialization map explicit.  Choose a compact \(T_f\)-orbit
\(K=T_f\cdot w_0\simeq T^2\) in the smooth stratum
\(W^\circ\simeq(\mathbb C^*)^2\).  On a relatively compact neighbourhood
of \(K\), average a Riemannian metric over \(T_f\) and horizontally lift
\(\partial_{\operatorname{Re}t}\) and
\(\partial_{\operatorname{Im}t}\).  If \(\xi_1^h,\xi_2^h\) are those lifts,
put
\[
 \xi_t=(\operatorname{Re}t)\xi_1^h+
       (\operatorname{Im}t)\xi_2^h.
\]
Choose a relatively compact \(T_f\)-invariant open neighbourhood \(K'\)
with \(K\Subset K'\).  For some \(r_0>0\), the time-one flow
\[
 \Theta:K'\times\Delta_{r_0}\longrightarrow N_0,
 \qquad \Theta(w,t)=\Phi^{\xi_t}_1(w),
\tag{5.9a}
\]
is a \(T_f\)-equivariant open embedding over \(\Delta_{r_0}\), with
\(\Theta(w,0)=w\).  Indeed \(df_0(\xi_t)=t\), so
\(f_0(\Theta(w,t))=t\); its differential is invertible along
\(K'\times\{0\}\), and injectivity follows first by applying \(f_0\) and
then from uniqueness for the time-one flow.  Shrinking \(K'\) and \(r_0\)
gives the claimed embedding.

For \(0<r<r_0\), set
\[
 A_r=f_0^{-1}(\Delta_r),\qquad
 Z_r=\Theta(K\times\Delta_r),\qquad
 K_t=\Theta(K\times\{t\}).
\]
The subset \(Z_r\) is closed in \(A_r\).  Indeed, if a sequence in \(Z_r\)
converges in \(A_r\), compactness of \(K\) gives a convergent subsequence of
its \(K\)-coordinates; its base coordinates converge to the image under
\(f_0\), which still lies in \(\Delta_r\), and the embedding \(\Theta\)
then identifies the limit with a point of \(Z_r\).

For \(0<|t|<r\), functoriality of relative cohomology gives the two
commutative squares
\[
\begin{CD}
H^3(A_r,A_r\setminus Z_r) @>>> H^3(A_r)\\
@V{u_0\;(\simeq)}VV                  @VV{\rho_0}V\\
H^3(W,W\setminus K) @>>> H^3(W)
\end{CD}
\qquad
\begin{CD}
H^3(A_r,A_r\setminus Z_r) @>>> H^3(A_r)\\
@V{u_t\;(\simeq)}VV                  @VV{\rho_t}V\\
H^3(F_t,F_t\setminus K_t) @>>> H^3(F_t).
\end{CD}
\tag{5.9b}
\]
The two left vertical arrows are isomorphisms.  By excision they are the
restrictions to the two slices of
\[
 H^3\bigl(K'\times\Delta_r,
          (K'\setminus K)\times\Delta_r\bigr),
\]
and homotopy invariance makes both restrictions isomorphisms.  Taking the
direct limit over shrinking \(r\), the maps \(\rho_0\) induce the
proper-base-change isomorphism
\[
 \varinjlim_{r\downarrow0}H^3(A_r;\mathbb Z)
 \xrightarrow{\ \sim\ }H^3(W;\mathbb Z).
\]
Choose the points \(t\) on one radial ray and identify the \(F_t\) with the
fixed marked fibre \(F\) by parallel transport.  Restriction to \(F_t\),
followed by this identification, is compatible under shrinking \(r\).
The nearby-cycle specialization is precisely the composite
\[
 \operatorname{sp}^3:
 H^3(W;\mathbb Z)
 \xleftarrow[\ \sim\ ]{\rho_0}
 \varinjlim_{r\downarrow0}H^3(A_r;\mathbb Z)
 \xrightarrow{\rho_{\rm near}}H^3(F;\mathbb Z);
\]
this is the neighbourhood realization of the unit
\(\mathbb Z_W\to\psi_{f_0}\mathbb Z\) in the comparison cited above.
Taking the direct limit in the two functorial squares therefore proves,
rather than assumes, that
\[
 H^3(W,W\setminus K)\longrightarrow H^3(W)
 \xrightarrow{\operatorname{sp}^3}H^3(F_t)
\tag{5.9c}
\]
is the transported supported class followed by the Gysin map for
\(K_t\hookrightarrow F_t\).

Under the marking of the torus fibre,
\[
 \Lambda_{\rm tor}=\mathbb Z\widehat w\oplus
                    \mathbb Z\widehat\delta
\]
is a saturated direct summand of \(\Lambda\), and \(K_t\) is a translate of
\((\Lambda_{\rm tor}\otimes\mathbb R)/\Lambda_{\rm tor}\).  Choose a
complementary summand \(\Lambda'\), and write
\(F_t\simeq T_f\times T_B\) as oriented real tori, where
\(T_B=(\Lambda'\otimes\mathbb R)/\Lambda'\), with the product orientation
chosen to agree with that of \(F_t\).  The real rank-two normal bundle of
the translate \(K_t\) is trivial and oriented.  The Thom isomorphism gives
\[
H^3(F_t,F_t\setminus K_t)\simeq H^1(T_f),
\]
and its map to absolute cohomology is
\[
 \iota_!(a)=a\otimes[T_B]^\vee
 \in H^1(T_f)\otimes H^2(T_B)\subset H^3(F_t).
\tag{5.9d}
\]
This is an injective primitive rank-two K\"unneth direct summand.  On the
central smooth stratum,
\(W^\circ\simeq(\mathbb C^*)^2\simeq T_f\times\mathbb R^2\) as an
oriented real manifold; hence the real rank-two normal bundle of \(K\) is
also trivial and oriented.  The same Thom calculation gives
\(H^3(W,W\setminus K)\simeq H^1(K)\simeq\mathbb Z^2\).
Hence (5.9c) has primitive rank-two image.  Since
\(\operatorname{rk}H^3(W)=2\), specialization in degree three is injective
and has precisely that primitive image.  In (5.9), this forces
\(\beta=\epsilon=0\), and consequently \(\delta'=1\).

We make the integral saturation argument explicit.  The unit
\(\mathbb Z_W\to\psi\) is a morphism from the complex concentrated in
cohomological degree zero to the Postnikov filtration of \(\psi\).  It
therefore induces a morphism of hypercohomology spectral sequences from the
single row
\[
 H^a(W;\mathbb Z)=E_2^{a,0}(\mathbb Z_W)
\]
to (5.8).  On the \(b=0\) row the induced \(E_2\)-map is the identity,
because \(\mathcal H^0\psi=\mathbb Z_W\).  Consequently, in total degree
\(r\), specialization factors through the last filtration step
\[
 H^r(W;\mathbb Z)\xrightarrow{\ \sim\ }
 E_\infty^{r,0}=F^rH^r(F;\mathbb Z)
 \lhook\joinrel\longrightarrow H^r(F;\mathbb Z)
 \tag{5.9e}
\]
whenever no differential
\(d_k:E_k^{r-k,k-1}\to E_k^{r,0}\) enters on any page.  For \(r=1\) this
is automatic; for \(r=2\) the only possible entering differential is
\(d_2^{0,1}\), and \(\alpha=0\).  Thus (5.9e) is an isomorphism onto the
specialization image in degrees one and two.  The quotient by \(F^r\) has
the finite filtration whose graded groups are the remaining
\(E_\infty^{a,r-a}\), \(a<r\).  In degree one these are kernels of maps of
free abelian groups.  In degree two they are
\(E_\infty^{1,1}=\mathbb Z\) and
\(E_\infty^{0,2}=\ker d_2^{0,2}\); here the possible subsequent map is
\(d_3^{0,2}\), and it is zero because its already determined rank
\(\epsilon\) is zero.  Both groups are therefore free.  Successive extensions of
free abelian groups split, so \(H^r(F;\mathbb Z)/F^r\) is free for
\(r=1,2\).  Hence the specialization images in those degrees are
primitive.

In degree three, (5.9c)--(5.9d) give a primitive rank-two lattice contained
in \(\operatorname{im}\operatorname{sp}^3\).  Since
\(H^3(W;\mathbb Z)\) itself has rank two, the specialization image has the
same rational span; a primitive lattice has no proper finite-index
overlattice in that span inside \(H^3(F;\mathbb Z)\).  Therefore the two
images coincide, proving both injectivity and primitivity in degree three.
In degree four, \(E_\infty^{4,0}=H^4(F;\mathbb Z)\), so specialization is an
isomorphism.  Thus specialization is injective with primitive image in every
degree.  Monodromy invariance gives
\[
 \operatorname{im}\operatorname{sp}^r
 \subseteq H^r(F;\mathbb Z)^{T_0}.
\]
By \cref{thm:cusp-lattices}, the invariant ranks are
\((1,2,4,2,1)\), exactly those in (5.7), and each invariant lattice is the
kernel of an integral endomorphism, hence primitive.  Two primitive
sublattices of the same rational subspace coincide: both equal the
intersection of that subspace with the ambient lattice.  This proves (5.1).
\end{proof}

\subsection{The global \texorpdfstring{\(E_2\)}{E2}-page}

The Leray spectral sequence used in this subsection is
\cite[Corollary~2.3.4]{Dimca2004}.
Let \(R^r=R^rf_*\mathbb Z\) and
\(\mathcal L^r=j^*R^r\).  The finite-fibre lattices have already been proved
as abstract lattice statements in \cref{prop:finite-indices}.  Their
identification with the stalks of the proposed family requires the explicit
radial-stalk compatibility below.  This compatibility is treated separately
in \cref{lem:finite-filling-stalk-map} and is not silently absorbed here.

\begin{proposition}[all Leray terms used below]
\label{prop:leray-e2}
\Verified\ The specialization theorem
\cref{thm:p0-integral-specialization}, the finite fillings of
\cref{prop:candidate-finite-fillings}, and
\cref{lem:finite-filling-stalk-map} give, for \(j=1,2\) and
\(0\le r\le4\), the finite-filling compatibility
\[
\tag{FF}
 N_j\ \text{strongly deformation retracts onto }S_j,\qquad
 H^r(S_j;\mathbb Z)\longrightarrow H^r(F;\mathbb Z)^{T_j}
 \ \text{is }\pi_j^*
\]
under the marked punctured-fibre identification.  Consequently,
\[
H^0(B,R^1)=\mathbb Z(12\gamma),\quad
H^0(B,R^2)=\mathbb Z(2q),\quad
H^0(B,R^3)=\mathbb Z(2\gamma uw),
\tag{5.10}
\]
\[
 H^1(B,R^1)=H^1(B,R^2)=0,
\tag{5.11}
\]
and, with \(\omega\) a generator of \(H^2(B;\mathbb Z)\),
\[
H^2(B,R^1)=\mathbb Z\omega[\delta],\qquad
H^2(B,R^2)=\mathbb Z\omega[\gamma\delta],\qquad
H^2(B,R^4)=\mathbb Z\omega[\operatorname{vol}].
\tag{5.12}
\]
These identifications are compatible with exterior products in \(V\).
\end{proposition}

\begin{proof}
Under the stated hypotheses, proper base change identifies the stalk of
\(R^r\) at \(p_0\) with \(H^r(W;\mathbb Z)\), while
\((j_*\mathcal L^r)_{p_0}=H^r(F;\mathbb Z)^{T_0}\).  The stalk map is
specialization and is an isomorphism by
\cref{thm:p0-integral-specialization}; hence its cokernel at \(p_0\) is zero.
At \(p_j\), condition \textup{(FF)} identifies the stalk with
\(H^r(S_j;\mathbb Z)\), the target with
\(H^r(F;\mathbb Z)^{T_j}\), and the stalk map with \(\pi_j^*\).  The
finite-fibre calculation \cref{prop:finite-indices} makes each of these maps
injective.  Away from the three special points the map is the identity of
the local system.  Stalkwise kernels and cokernels therefore give the exact
sequence
\[
0\longrightarrow R^r\longrightarrow j_*\mathcal L^r
\longrightarrow Q^r\longrightarrow0,
\tag{5.13}
\]
with
\[
Q^1=(\mathbb Z/3)_{p_1}\oplus(\mathbb Z/4)_{p_2},
\quad Q^2=(\mathbb Z/2)_{p_2},
\quad Q^3=(\mathbb Z/2)_{p_2},
\quad Q^4=(\mathbb Z/3)_{p_1}\oplus(\mathbb Z/4)_{p_2}.
\]
The degree-one quotients have orders three and four and are generated by the
class of \(\gamma\).  In degree two the \(p_1\)-quotient is zero, while
\[
 (\textstyle\bigwedge^2V)^{T_2}/
 \{a\gamma\eta_2+bq:a\equiv b\pmod2\}\simeq\mathbb Z/2
\]
is generated by the class of \(q\).  In degree three the \(p_1\)-quotient
is zero and the \(p_2\)-quotient has order two, generated by the class of
\(\gamma uw\).  This last identification uses the actual twist
\(v_2=-\varepsilon'\), for which the source's admissibility quantity
\(\eta_2(v_2)\) is zero and hence divisible by four.  In degree four,
pullback multiplies the orientation class by the covering
degree \(m_j\), giving the displayed \(Q^4\).  Hence the skyscraper sheaves
have exactly the asserted stalks; no claim for an arbitrary twist is being
imported.
By \cref{thm:coinvariants},
\[
V^G=\mathbb Z\gamma,\quad
(\textstyle\bigwedge^2V)^G=\mathbb Zq,\quad
(\textstyle\bigwedge^3V)^G=\mathbb Z\gamma uw,
\tag{5.14}
\]
and
\[
V_G=\mathbb Z[\delta],\qquad
(\textstyle\bigwedge^2V)_G=\mathbb Z[\gamma\delta],\qquad
[q]=12[\gamma\delta].
\tag{5.15}
\]
Since \(\det T_1=\det T_2=1\), the action on
\(\bigwedge^4V=\mathbb Z\operatorname{vol}\) is trivial and
\[
 (\textstyle\bigwedge^4V)_G=\mathbb Z[\operatorname{vol}].
\tag{5.15a}
\]
The maps from the invariant groups in (5.14) onto the corresponding
skyscraper fibres are surjective, with kernels
\(12\mathbb Z\gamma\), \(2\mathbb Zq\), and
\(2\mathbb Z\gamma uw\).  This proves (5.10).

Because each \(Q^r\) is a skyscraper, it has no positive cohomology; the
long exact sequence of (5.13) therefore identifies
\(H^2(B,R^r)\) with \(H^2(B,j_*\mathcal L^r)\).
For any free local system \(\mathcal L\) on \(B^\circ\), with stalk \(M\),
the quotient \(j_*\mathcal L/j_!\mathcal L\) is a skyscraper.  Hence
\[
 H^2(B,j_*\mathcal L)\simeq H_c^2(B^\circ,\mathcal L)\simeq
 H_0(B^\circ,\mathcal L)=M_G,
\tag{5.16}
\]
where the middle isomorphism is Poincar\'e duality with local coefficients
\cite[\S3.H]{Hatcher2002}; see also
\href{https://doi.org/10.1007/978-1-4612-0647-7}{G.~E.~Bredon,
\emph{Sheaf Theory}, 2nd ed., Chapter~VI}.
Cap product with the locally finite fundamental class makes (5.16) natural
for coefficient pairings, proving (5.12) and the product assertion.

It remains to prove (5.11), for which we spell out the parabolic-cohomology
quotient.  The low-degree Leray sequence for \(j:B^\circ\hookrightarrow B\)
contains
\[
 0\longrightarrow H^1(B,j_*\mathcal L)
 \longrightarrow H^1(B^\circ,\mathcal L)
 \longrightarrow H^0(B,R^1j_*\mathcal L).
\]
If \(I_i=\langle t_i\rangle\) is the local inertia group at \(p_i\), then
\[
 (R^1j_*\mathcal L)_{p_i}
 =H^1(I_i,M)=M/(T_i-I)M.
\]
Consequently the last arrow sends the class of a cocycle \(c\) to the three
classes of \(c(t_i)\), and \(H^1(B,j_*\mathcal L)\) is exactly the kernel
defined by the three parabolic conditions below.  Use left cocycles, so
\(c(gh)=c(g)+g c(h)\).  A cocycle on
\(\pi_1(B^\circ)=F(t_1,t_2)\) is determined by
\(x=c(t_1)\) and \(y=c(t_2)\).  Extension across the three punctures is the
parabolic condition
\[
 x\in(T_1-I)M,\qquad y\in(T_2-I)M,\qquad
 c(t_0)\in(T_0-I)M,
 \qquad t_0=(t_1t_2)^{-1}.
\]
Write \(y=(T_2-I)m\) and subtract the coboundary of \(m\); this normalizes
\(y=0\).  A further coboundary preserves the normalization precisely when
its defining element lies in \(M^{T_2}\), and then changes \(x\) by an
element of \((T_1-I)M^{T_2}\).  With \(y=0\),
\(c(t_0)=-T_0x\).  Since \(T_0\) preserves
\((T_0-I)M\), the last parabolic condition is equivalent to
\(x\in(T_0-I)M\).  Therefore
\[
 H^1(B,j_*\mathcal L)\simeq
 \frac{(T_1-I)M\cap(T_0-I)M}{(T_1-I)M^{T_2}}.
\tag{5.17}
\]
For \(M=V\),
\[
(T_1-I)V=\langle2\gamma+u+w,2\gamma-u\rangle,\qquad
(T_0-I)V=\langle\gamma,u\rangle,
\]
and their intersection is
\(\mathbb Z(2\gamma-u)=(T_1-I)\mathbb Z\eta_2\).  For
\(M=\bigwedge^2V\),
\[
\begin{aligned}
(\textstyle\bigwedge^2T_1-I)\textstyle\bigwedge^2V
=\langle&\gamma u,\gamma w,2u\delta+w\delta,\\
&uw+u\delta-w\delta-6\gamma\delta\rangle,
\end{aligned}
\]
whereas
\[
(\textstyle\bigwedge^2T_0-I)\textstyle\bigwedge^2V
=\langle\gamma u,\gamma w+u\delta\rangle.
\]
Their intersection is \(\mathbb Z\gamma u\), equal to
\((\bigwedge^2T_1-I)\mathbb Z(\gamma\eta_2)\).  Thus both groups in
(5.17) vanish.  The maps
\(H^0(B,j_*\mathcal L^r)\to H^0(B,Q^r)\) were proved surjective above for
\(r=1,2\), and skyscraper sheaves have no positive cohomology.  The long
exact sequence of (5.13) therefore gives (5.11) for \(R^1,R^2\).
\end{proof}

\subsection{The three transgressions}

\begin{lemma}[the primitive first transgression]
\label{lem:first-transgression}
\Verified\ For the constructed map and the based meridians of
\cref{prop:geometric-pi1}, one may orient \(\omega\) so that
\[
 d_2^{0,1}(12\gamma)=\omega.
\tag{5.18}
\]
\end{lemma}

\begin{proof}
Choose disjoint discs \(D_i\) about the special points and a pair-of-pants
neighbourhood \(V_0\) of their complement.  The class
\(\theta=12\gamma\) is a section of \(R^1f_*\mathbb Z\) over each \(D_i\):
at \(p_0\) this is \cref{thm:p0-integral-specialization}, and at \(p_j\)
it is the degree-one finite-fibre lattice.  It is also a section over
\(V_0\), because it is globally monodromy invariant.  For every member
\(U\in\{D_0,D_1,D_2,V_0\}\), the local Leray five-term sequence contains
\[
 H^1(f^{-1}(U);\mathbb Z)\longrightarrow
 H^0(U,R^1f_*\mathbb Z)\longrightarrow
 H^2(U,R^0f_*\mathbb Z)=H^2(U;\mathbb Z)=0.
\]
Thus the section \(\theta\) has a cohomological lift
\(\Theta_U\in H^1(f^{-1}(U);\mathbb Z)\) on every member of the cover.
On an overlap \(D_i\cap V_0\), the two lifts have the same image in
\(H^0(R^1f_*\mathbb Z)\); local five-term exactness therefore identifies
their difference with a class in
\(H^1(D_i\cap V_0,R^0f_*\mathbb Z)=H^1(D_i\cap V_0;\mathbb Z)\).
These three differences form the \v{C}ech connecting cocycle used below.

Use clockwise meridians.  The affine deck relation at \(p_j\) is
\[
 \widehat\rho_j^{m_j}=t_{v_j}.
\tag{5.19}
\]
Put \(\Theta_j=\Theta_{D_j}\).  Since \(\Theta_j\) is integral and restricts
to \(\theta\) on the fibre, (5.19) gives
\(m_j\Theta_j(\widehat\rho_j)=\Theta_j(t_{v_j})=\theta(v_j)\).  Therefore
\[
 \Theta_1(\widehat\rho_1)=12\gamma(v_1)/3=4,
 \qquad
 \Theta_2(\widehat\rho_2)=12\gamma(v_2)/4=-3.
\tag{5.20}
\]
By \cref{prop:geometric-pi1}, the toric meridian
\(\widetilde a_0\) is the
boundary of the holomorphic disc
\(\{(1,1,t):|t|\le\epsilon\}\).  It is therefore null-homotopic and the
local lift evaluates to zero.

Write \(\alpha_i=\Theta_{V_0}(\widehat\rho_i)\).  By the exact based cusp
relation (5.23a), the three annular discrepancies have
meridian values
\[
 4-\alpha_1,\qquad -3-\alpha_2,\qquad \alpha_1+\alpha_2,
\tag{5.21}
\]
whose sum is one.  Mayer--Vietoris on the base identifies
\[
H^1\!\left(\coprod_i(D_i\cap V_0)\right)/\operatorname{im}H^1(V_0)
\simeq H^2(B;\mathbb Z).
\]
Under meridian evaluation the numerator is \(\mathbb Z^3\), the image is
the sum-zero lattice, and the connecting map is the signed sum of the three
coordinates.

We now identify this connecting class with the Leray differential rather
than taking that identification as an extra hypothesis.  Put
\(K=Rf_*\mathbb Z_X\).  The canonical truncation triangle
\cite[Tag~08J5]{StacksProject} is
\[
 R^0f_*\mathbb Z_X\longrightarrow\tau_{\le1}K
 \longrightarrow R^1f_*\mathbb Z_X[-1]
 \xrightarrow{\ \kappa\ }R^0f_*\mathbb Z_X[1].
\tag{5.21a}
\]
After applying \(R\Gamma(B,-)\), the connecting homomorphism
\[
 H^0(B,R^1f_*\mathbb Z)\xrightarrow{\,H^0(\kappa)\,}
 H^2(B,R^0f_*\mathbb Z)
\tag{5.21b}
\]
is precisely \(d_2^{0,1}\) in the hypercohomology, hence Leray, spectral
sequence; this is the canonical-truncation construction of that spectral
sequence \cite[Tag~0FLJ]{StacksProject}.  On the chosen disc--pair-of-pants
cover, functorial Mayer--Vietoris for the triangle
\cite[Tag~0F16]{StacksProject} computes \(H^0(\kappa)\) by choosing the local
degree-one lifts \(\Theta_U\) and taking their differences on the annular
overlaps.  Those differences are exactly the cocycle (5.21).  Thus the
signed sum just computed is \(d_2^{0,1}(\theta)\).
Equation (5.21) therefore represents a primitive generator; orient \(\omega\)
to obtain (5.18).
\end{proof}

\begin{lemma}[the higher transgressions]
\label{lem:higher-transgressions}
\Verified\ In the integral Leray spectral sequence of the constructed map,
\[
 d_2^{0,2}(2q)=\omega[\delta],
 \qquad
 d_2^{0,3}(2\gamma uw)=\omega[\gamma\delta].
\tag{5.22}
\]
\end{lemma}

\begin{proof}
Use the multiplicative Leray spectral sequence obtained from a filtered
multiplicative cochain resolution.  Its \(E_2\)-product is the natural cup
product
\[
 H^p(B,R^q)\otimes H^{p'}(B,R^{q'})
 \longrightarrow H^{p+p'}(B,R^{q+q'}),
\]
and every \(d_r\) is a graded derivation.  This is the standard filtered-DGA
construction; compare the filtered-complex spectral sequence
\cite[Tags~012K and~0FLJ]{StacksProject} with the derived cup product
\cite[\S20.31]{StacksProject}.  The generator identifications and their
products are natural by the last sentence of
\cref{prop:leray-e2}.
Put \(\theta=12\gamma\), \(a=2q\), and \(b=2\gamma uw\).  By
\cref{prop:leray-e2}, write
\[
 d_2a=r\omega[\delta],\qquad d_2b=s\omega[\gamma\delta].
\]
In the \(E_2\)-ring,
\[
 \theta a=12b,\qquad ab=0,\qquad [q]=12[\gamma\delta].
\]
The derivation rule and \(|\theta|=1\) give
\[
12d_2b=d_2(\theta a)=d_2\theta\cdot a-\theta\cdot d_2a,
\]
hence \(s=2-r\).  Since \(|a|=2\),
\(0=d_2(ab)=d_2a\cdot b+a\cdot d_2b\).  The exterior products are
\[
 \delta\wedge2\gamma uw=-2\operatorname{vol},
 \qquad
 2q\wedge\gamma\delta=2\operatorname{vol}.
\]
Therefore \(-2r+2s=0\), so \(r=s\).  Together with \(s=2-r\), this gives
\(r=s=1\), proving (5.22).
\end{proof}

\begin{theorem}[integral cohomology in low degrees]
\label{thm:h2-vanishing}
\Verified\ For the constructed holomorphic map,
\[
 H^1(X;\mathbb Z)=H^2(X;\mathbb Z)=H^3(X;\mathbb Z)=0,
 \qquad H_2(X;\mathbb Z)=0.
\tag{5.23}
\]
\end{theorem}

\begin{proof}
The Leray terms in total degree two are
\[
E_2^{2,0}=\mathbb Z\omega,\qquad E_2^{1,1}=0,\qquad
E_2^{0,2}=\mathbb Z(2q).
\]
By \cref{lem:first-transgression,lem:higher-transgressions}, the incoming
map to the first term and the outgoing map from the last term are
isomorphisms.  Thus all three \(E_\infty\)-pieces vanish and \(H^2(X)=0\).
The first map has zero kernel, giving \(H^1(X)=0\).  In total degree three,
the quotient of \(E_2^{2,1}\) by the image of \(E_2^{0,2}\) is zero,
\(E_2^{1,2}=0\), and the outgoing map from \(E_2^{0,3}\) is injective by
(5.22).  Hence \(H^3(X)=0\).  No later differential can meet these terms
because the base has cohomological dimension two.

The universal coefficient sequences show both
\(\operatorname{Hom}(H_2(X),\mathbb Z)=0\) and
\(\operatorname{Ext}(H_2(X),\mathbb Z)=0\).  Since \(H_2(X)\) is finitely
generated, it has neither a free nor a torsion summand, so it is zero.
\end{proof}

\subsection{Based van Kampen and sphere recognition}

\begin{proposition}[the geometric fundamental-group presentation]
\label{prop:geometric-pi1}
\Verified\ For the punctured-overlap gluings of
\cref{prop:candidate-cusp-overlap,prop:candidate-finite-fillings}, the based
van Kampen presentation is
\[
 \pi_1(X)\simeq
 \langle c,x,y\mid c\text{ central},\ xy=1,\ x^3=c,\ y^4=c^{-1}\rangle
 \simeq1.
\tag{5.24}
\]
\end{proposition}

\begin{proof}
Choose \(b_*\) in a pair-of-pants core of \(B^\circ\), put \(x_*\) on the
holomorphic zero section above it, and choose three nonintersecting base
whiskers from \(b_*\) to clockwise boundary circles.  Lift all whiskers by
the zero section.  With the ordering compatible with \(g_1g_2g_0=1\), the
resulting based loops obey
\[
 \widehat\rho_1\widehat\rho_2\widetilde a_0=1.
\tag{5.23a}
\]
The period-family action is linear and hence preserves the zero section.
It splits the homotopy sequence over the pair of pants, giving
\[
 \pi_1(f^{-1}(B^\circ))=
 \Lambda\rtimes F(\widehat\rho_1,\widehat\rho_2),
\]
with conjugation actions \(A_1,A_2\).  The clockwise meridian at
\(p_j\) lifts to the affine generator in (2.9), with its exact
sign and based translation.  To verify this in the actual
real marking, put
\[
 \alpha(s)=\frac{\log s_j}{2\pi i},\qquad
 J_s=\Pi(s)^{-1}i\Pi(s),\qquad
 a_j(s)=\operatorname{Re}\alpha(s)\,v_j+
        \operatorname{Im}\alpha(s)\,J_s v_j.
\]
The overlap translation by
\(\alpha(s)\Pi(s)v_j\) is translation by \(a_j(s)\) in the real
marking.  Thus the inverse image of the period-family zero section
is \(x(s)=-a_j(s)\).

The covariance of \(J_s\) is a consequence of the proved period
identity (2.8e), not an additional assumption.  That identity is
\[
 R_{g_j}(s)\Pi(s)=\Pi(g_js)A_j.
\]
The map \(R_{g_j}(s)\) is complex linear.  Hence
\[
 \begin{split}
 J_{g_js}A_j
 &=\Pi(g_js)^{-1}i\,\Pi(g_js)A_j\\
 &=\Pi(g_js)^{-1}i\,R_{g_j}(s)\Pi(s)\\
 &=\Pi(g_js)^{-1}R_{g_j}(s)i\,\Pi(s)
 =A_jJ_s.
 \end{split}
\]
The original vectors are
\[
 v_1=(1,2,-4,0)^{\mathsf t},\qquad
 v_2=(-1,-3,3,0)^{\mathsf t},
\]
in the basis
\((\widehat\gamma,\widehat u,\widehat w,\widehat\delta)\).
Multiplication by the printed matrices gives
\[
 A_1v_1=(1,\,6-4,\,-6-2+4,\,-2+2)^{\mathsf t}
        =(1,2,-4,0)^{\mathsf t},
\]
\[
 A_2v_2=(-1,\,-3,\,6-3,\,-3+3)^{\mathsf t}
        =(-1,-3,3,0)^{\mathsf t}.
\]
Thus \(A_jv_j=v_j\) exactly.  Clockwise continuation by \(g_j\)
changes the chosen logarithm to
\(\alpha(g_js)=\alpha(s)-1/m_j\).  Using this real change, the
unchanged imaginary part, and \(J_{g_js}A_j=A_jJ_s\), one obtains
\[
 a_j(g_js)=A_ja_j(s)-v_j/m_j.
\]
Therefore the affine generator sends the initial lift of the
zero section to its actual continued endpoint:
\[
 A_j(-a_j(s))+v_j/m_j=-a_j(g_js).
\]
Changing the logarithm branch by an integer changes the initial
lift by an integer multiple of \(v_j\).  Such a translation
commutes with the affine generator because \(A_jv_j=v_j\), so the
endpoint identification does not introduce a translated conjugate.
The fixed zero-section whiskers retain the based generator.
After \(m_j\) clockwise iterations, the coordinate returns to the
same point and \(\alpha\) has decreased by one.  Thus the endpoint
of \(-a_j\) has increased by \(v_j\), in agreement with (2.8f).
Consequently
\[
 \widehat\rho_j^{m_j}=t_{v_j}.
 \tag{5.23b}
\]
The radial model of \cref{lem:finite-filling-stalk-map} shows that attaching
\(N_j\) imposes exactly this relation.  Since the global and local whiskers
are the same zero-section path under the overlap, conjugating back to
\(x_*\) adds no fibre drift.

For completeness, the cusp covering has no hidden fundamental group.  A loop
in the infinite toric variety has compact image and therefore lies in the
open \(X_{\Sigma_{\rm tor}}\) associated with a finite subfan
\(\Sigma_{\rm tor}\).  Enlarge that subfan by the rays through
\((0,1),(e_1,1),(e_2,1)\).  Write
\[
 N_{\Sigma_{\rm tor}}
 :=\bigl\langle|\Sigma_{\rm tor}|\cap N'\bigr\rangle_{\mathbb Z}.
\]
The three displayed primitive generators span \(N'\), so the classical
formula
\(\pi_1(X_{\Sigma_{\rm tor}})=N'/N_{\Sigma_{\rm tor}}\)
\cite[Theorem~12.1.10]{CoxLittleSchenck2011} makes the resulting finite toric
open simply connected; the original loop therefore has a null-homotopy in
that open.  Thus \(Y\) is simply connected.  The one-parameter subgroup in
the third torus coordinate acts by \(t\mapsto r t\), \(r>0\).  Given the
compact image of a null-homotopy, choose \(r\) so small that its translate
lies in \(|t|<\epsilon\); the path \(r_s=(1-s)+sr\) gives a homotopy from the
original loop to its translate inside \(Y_\epsilon\).  Hence
\(Y_\epsilon\) is also simply connected.  The verified free quotient
\(N_0=Y_\epsilon/\bar\Lambda\) therefore has universal cover
\(Y_\epsilon\), so
\(\pi_1(N_0)=\bar\Lambda\).  Equation (2.8iii) identifies a fibre-lattice
loop \(\lambda\in\Lambda\) with the deck transformation
\(\bar\lambda\); hence the inclusion of a nearby fibre is exactly
\[
 \Lambda\longrightarrow\bar\Lambda,\qquad
 \lambda\longmapsto\bar\lambda,
\tag{5.23c}
\]
with kernel
\(\Lambda_{\rm tor}=\langle\widehat w,\widehat\delta\rangle\).
On the zero section the same overlap sends \((z,0)\) to
\((1,1,t_c)\).  In the toric chart of the triangle
\(\{0,e_1,e_2\}\), whose coordinates may be written
\((t_c/(x_1x_2),x_1,x_2)\), this is the holomorphic disc
\(t_c\mapsto(t_c,1,1)\).  Thus the cusp meridian
\(\widetilde a_0\) is null-homotopic, with no fibre translation.  Combining
this fact with (5.23a), the cusp attachment imposes precisely
\(\widehat\rho_1\widehat\rho_2=1\).

Van Kampen now gives
\[
\left\langle\Lambda,\widehat\rho_1,\widehat\rho_2\ \middle|\
\widehat\rho_j\lambda\widehat\rho_j^{-1}=A_j\lambda,
\widehat\rho_1^3=t_{v_1},\widehat\rho_2^4=t_{v_2},
\widehat\rho_1\widehat\rho_2=1,
\Lambda_{\rm tor}=1\right\rangle.
\]
The normal closure of \(\Lambda_{\rm tor}\) under \(A_1,A_2\) is
\(\ker\gamma=\langle\widehat u,\widehat w,\widehat\delta\rangle\); for
example \(A_1\widehat w=\widehat u-\widehat w\).  The image of \(\Lambda\)
is therefore generated by \(c=[\widehat\gamma]\), which is central because
\(\gamma\) is invariant.  Since \(\gamma(v_1)=1\) and
\(\gamma(v_2)=-1\), the presentation becomes (5.24).  It is the case
\((\ell_0,\ell_1,\ell_2)=(0,1,-1)\) of
\cref{prop:terminal-group}, where \(p=-1\); hence the group is trivial.
\end{proof}

\begin{corollary}[topology and the third Chern number]
\label{cor:candidate-sphere}
\Verified\ The underlying smooth manifold of the constructed threefold is
diffeomorphic to the standard six-sphere, and
\[
 \langle c_3(T_X),[X]\rangle=2.
\tag{5.25}
\]
\end{corollary}

\begin{proof}
Combine \cref{thm:h2-vanishing,prop:geometric-pi1}.  The universal
coefficient theorem and Poincar\'e duality give the integral homology of
\(S^6\), exactly as in \cref{lem:cohomological-closure}.  The complete
recognition argument of \cref{thm:sphere-recognition} then gives
\(X\cong_{\rm diff}S^6\).  Its Euler characteristic is two.  On a complex
threefold the top Chern class of the complex tangent bundle is the Euler
class of the underlying oriented real tangent bundle, proving (5.25).
\end{proof}

\paragraph{Independent role of the collapse route.}
The proof above is independent of the cellular route.  The candidate's cusp
collapse, its integral homology map, and its equality with nearby-cycle
specialization are now independently verified in
\cref{prop:verified-cellular-collapse,prop:verified-collapse-homology,%
thm:collapse-equals-specialization}.  The actual finite-band and cusp
Mayer--Vietoris maps, their primitive blocks, and their commuting comparison
with the differentials proved in this section are now
\cref{prop:actual-finite-band-maps,prop:cusp-mv-blocks,%
thm:mv-leray-comparison}.  This supplies a second complete integral topology
certificate; it remains logically independent of the sphere-recognition
conclusion already obtained in this section.
